
\begingroup
\setlength{\LTleft}{0pt}
\setlength{\LTright}{0pt}
\setlength{\tabcolsep}{3pt}
\renewcommand{\arraystretch}{1.12}
\newcommand{\SplitAuditName}[1]{\textit{\hspace{0pt}#1}}
\begin{longtable}{@{}>{\raggedright\arraybackslash}p{0.22\linewidth} >{\raggedright\arraybackslash}p{0.15\linewidth} >{\raggedright\arraybackslash}p{0.115\linewidth} >{\raggedright\arraybackslash}p{\dimexpr0.515\linewidth-6\tabcolsep\relax}@{}}
\caption{UCR split-audit categorization (128 datasets): Concern 32.03\%; Unclear 50.78\%; Separated 5.47\%; Related periods 5.47\%; Generated 6.25\%.}\label{tab:splits} \\
\toprule
Dataset & Relevant new case & Verdict & What the split tests and why a concern can rise \\
\midrule
\endfirsthead
\multicolumn{4}{l}{\emph{UCR split audit (continued)}} \\
\toprule
Dataset & Relevant new case & Verdict & What the split tests and why a concern remains \\
\midrule
\endhead
\midrule
\multicolumn{4}{r}{\emph{continues on the next page}} \\
\endfoot
\bottomrule
\endlastfoot
\hypertarget{split-audit:ACSF1}{}\SplitAuditName{ACSF1} & New appliances & Unclear & The 100/100 split matches the source's Test Protocol 1.0, which trains on the first measurement session and tests on the second session of the same appliances. The archive does not say which protocol was used, and the series length does not match the documented one-hour sessions. Whether the test presents familiar appliances again or appliances the model has never seen therefore cannot be established. A concern arises only if the result is presented as performance on new appliances. \\
\hypertarget{split-audit:Adiac}{}\SplitAuditName{Adiac} & New independently collected samples & Unclear & The cases are outlines of tiny biological specimens. We do not know whether samples grown from the same biological source, or photographed on the same microscope slide, occur in both sets. They could share growth or imaging conditions. The records do not let us check whether the test samples were collected independently. \\
\hypertarget{split-audit:AllGestureWiimoteX}{}\SplitAuditName{AllGestureWiimoteX} & New people & Unclear & The data contain gestures from ten people. The archive shuffled and sampled recordings, but does not say that all recordings from a person were kept in one set. We therefore do not know whether test users are new. This matters if the intended application is to recognise gestures from people the model has never seen. \\
\hypertarget{split-audit:AllGestureWiimoteY}{}\SplitAuditName{AllGestureWiimoteY} & New people & Unclear & Same as \hyperlink{split-audit:AllGestureWiimoteX}{\SplitAuditName{AllGestureWiimoteX}}. \\
\hypertarget{split-audit:AllGestureWiimoteZ}{}\SplitAuditName{AllGestureWiimoteZ} & New people & Unclear & Same as \hyperlink{split-audit:AllGestureWiimoteX}{\SplitAuditName{AllGestureWiimoteX}}. \\
\hypertarget{split-audit:ArrowHead}{}\SplitAuditName{ArrowHead} & New physical objects & Unclear & Each case is an outline of an arrowhead. The source assigned outlines to a 36/175 split at random. We cannot check whether multiple views or representations of the same object were kept together. \\
\hypertarget{split-audit:BME}{}\SplitAuditName{BME} & New generated examples & Generated & This is a generated task rather than measurements from people, devices or specimens. The corresponding source-separation question is therefore not assessed here. Its split would instead need to be examined through the generation procedure; synthetic origin alone is not evidence of leakage. \\
\hypertarget{split-audit:Beef}{}\SplitAuditName{Beef} & New physical samples & Unclear & The data are spectral measurements of beef samples. It is not documented whether repeated measurements of one physical sample were kept in the same set. Without that information, performance on new samples cannot be distinguished from performance on further measurements of familiar samples. \\
\hypertarget{split-audit:BeetleFly}{}\SplitAuditName{BeetleFly} & New source images & Unclear & The cases are outlines extracted from images of beetles and flies. We do not know how the original images were assigned to training and testing. If different versions of one image occur in both sets, the model could recognise a familiar shape. Such shared images have not been verified. \\
\hypertarget{split-audit:BirdChicken}{}\SplitAuditName{BirdChicken} & New source images & Unclear & Same as \hyperlink{split-audit:BeetleFly}{\SplitAuditName{BeetleFly}}. \\
\hypertarget{split-audit:CBF}{}\SplitAuditName{CBF} & New generated examples & Generated & This is a generated task rather than measurements from people, devices or specimens. The corresponding source-separation question is therefore not assessed here. Its split would instead need to be examined through the generation procedure; synthetic origin alone is not evidence of leakage. \\
\hypertarget{split-audit:Car}{}\SplitAuditName{Car} & New vehicles or recordings & Unclear & The outlines were extracted from traffic videos. Several frames can show the same vehicle under very similar conditions. The documentation does not establish whether all frames of a vehicle or video were kept in one set, so performance on new vehicles or capture occasions cannot be checked. \\
\hypertarget{split-audit:Chinatown}{}\SplitAuditName{Chinatown} & New days or later periods & Related periods & A test day may be close to a training day and share its season or local events. These are different days, so this is not automatically leakage. No source says whether the split holds out a later period that follows all the training days. \\
\hypertarget{split-audit:ChlorineConcentration}{}\SplitAuditName{ChlorineConcentration} & New simulation runs & Concern & The audit describes cases taken at nearby times in one simulated system. A training case and test case can therefore describe almost the same system conditions. The result tests more cases from that simulation. It does not show whether the model works on a new simulation run; similarity between cases would need a separate check. \\
\hypertarget{split-audit:CinCECGTorso}{}\SplitAuditName{CinCECGTorso} & New recording occasions & Concern & The task is to identify four people from their ECG measurements. The archive defines each class as one person, so training and test share those people, and very likely one recording each. The same identities are expected in an identification task, but new recordings are not tested. The model may recognise characteristics of the recording as well as the person. \\
\hypertarget{split-audit:Coffee}{}\SplitAuditName{Coffee} & New physical samples & Unclear & The cases are spectral measurements of coffee. We do not know whether a sample was measured more than once and whether those measurements were kept together. If one measurement is in training and another in testing, the test sample is familiar. The missing sample information prevents us from checking this possibility. \\
\hypertarget{split-audit:Computers}{}\SplitAuditName{Computers} & New appliances or households & Unclear & The source measured each physical appliance over many days. We do not know whether all days from an appliance or household were kept in one set. If the same appliance appears in both sets, the model may recognise its usual pattern. That would be a narrower test than recognising appliance types in new homes. \\
\hypertarget{split-audit:CricketX}{}\SplitAuditName{CricketX} & New people & Unclear & The data contain repeated gestures, but the files do not say which person produced each case. We cannot check whether all gestures from a person were kept in one set. A test gesture might come from someone already represented in training. Performance on people the model has never seen is therefore not established. \\
\hypertarget{split-audit:CricketY}{}\SplitAuditName{CricketY} & New people & Unclear & Same as \hyperlink{split-audit:CricketX}{\SplitAuditName{CricketX}}. \\
\hypertarget{split-audit:CricketZ}{}\SplitAuditName{CricketZ} & New people & Unclear & Same as \hyperlink{split-audit:CricketX}{\SplitAuditName{CricketX}}. \\
\hypertarget{split-audit:Crop}{}\SplitAuditName{Crop} & New fields, scenes or seasons & Concern & The cases are pixel sequences from one scene and season. Nearby pixels may describe the same field and growing conditions. The value check also finds 2,753 of 16,800 test cases identical to training cases (Table~\ref{tab:duplicates}). This makes the test different from evaluating new fields, scenes or seasons. \\
\hypertarget{split-audit:DiatomSizeReduction}{}\SplitAuditName{DiatomSizeReduction} & New biological sources & Unclear & The sources explain size variation within each class by successive generations of a clone, but do not say how the split was drawn or whether training and test come from the same clone. If they do, a new specimen from a familiar line is a narrower test than a specimen from an independently collected source. The intended application should make clear which of these is relevant. \\
\hypertarget{split-audit:DistalPhalanxOutlineAgeGroup}{}\SplitAuditName{DistalPhalanxOutlineAgeGroup} & New patients & Unclear & The cases are finger-bone outlines extracted from X-rays. We do not know whether one patient contributed several images or whether those images were kept together. If one patient appears in both sets, the model may benefit from familiar anatomy. Patient identifiers are needed to check performance on new patients. \\
\hypertarget{split-audit:DistalPhalanxOutlineCorrect}{}\SplitAuditName{DistalPhalanxOutlineCorrect} & New patients & Unclear & Same as \hyperlink{split-audit:DistalPhalanxOutlineAgeGroup}{\SplitAuditName{DistalPhalanxOutlineAgeGroup}}. \\
\hypertarget{split-audit:DistalPhalanxTW}{}\SplitAuditName{DistalPhalanxTW} & New patients & Unclear & Same as \hyperlink{split-audit:DistalPhalanxOutlineAgeGroup}{\SplitAuditName{DistalPhalanxOutlineAgeGroup}}. \\
\hypertarget{split-audit:DodgerLoopDay}{}\SplitAuditName{DodgerLoopDay} & New days or later periods & Related periods & Days from one traffic detector were assigned at random. Training and test days can be part of the same baseball event or share similar traffic conditions. The test contains different days, but does not specifically evaluate a later period or another detector. This relationship should be considered when interpreting the score. \\
\hypertarget{split-audit:DodgerLoopGame}{}\SplitAuditName{DodgerLoopGame} & New days or later periods & Related periods & Same as \hyperlink{split-audit:DodgerLoopDay}{\SplitAuditName{DodgerLoopDay}}. \\
\hypertarget{split-audit:DodgerLoopWeekend}{}\SplitAuditName{DodgerLoopWeekend} & New days or later periods & Related periods & Same as \hyperlink{split-audit:DodgerLoopDay}{\SplitAuditName{DodgerLoopDay}}. \\
\hypertarget{split-audit:ECG200}{}\SplitAuditName{ECG200} & New recording sessions of a known patient & Concern & Each case is a heartbeat. The source thesis drew all 200 beats from one Holter recording of a single patient, so training and test beats share the patient and the recording session. For patient-specific monitoring a new case is a new recording session of that patient, which the split cannot test. \\
\hypertarget{split-audit:ECG5000}{}\SplitAuditName{ECG5000} & New patients or recording occasions & Concern & The source record identified by the audit contains heartbeats from one patient during one long recording. Training and test beats therefore share that patient and recording. This can test more beats from that record, but it cannot establish performance on another patient or after the patient is recorded again. \\
\hypertarget{split-audit:ECGFiveDays}{}\SplitAuditName{ECGFiveDays} & Independent recordings & Concern & The label says which of two recording dates produced a heartbeat. Training and test beats come from those same two recordings of one person. The model may distinguish conditions specific to those recordings. The test does not show that the distinction holds when further independent recordings are collected. \\
\hypertarget{split-audit:EOGHorizontalSignal}{}\SplitAuditName{EOGHorizontalSignal} & New people or recording occasions & Unclear & The eye-movement samples come from six people. The audit describes splitting individual movements rather than keeping all movements from a person together. We cannot establish performance on unfamiliar users. Sharing users can suit a personalised system, but its test should explain whether they were measured again on a separate occasion. \\
\hypertarget{split-audit:EOGVerticalSignal}{}\SplitAuditName{EOGVerticalSignal} & New people or sessions & Unclear & Same as \hyperlink{split-audit:EOGHorizontalSignal}{\SplitAuditName{EOGHorizontalSignal}}. \\
\hypertarget{split-audit:Earthquakes}{}\SplitAuditName{Earthquakes} & New periods or event groups & Related periods & The time windows are described as non-overlapping, so shared time samples are not an established problem here. However, windows from one earthquake sequence can be assigned to both sets. They may share the same broader event. The split does not specifically test later periods or separately collected event groups. \\
\hypertarget{split-audit:ElectricDevices}{}\SplitAuditName{ElectricDevices} & New appliances or households & Concern & The source kept each household in one set, but the shipped files differ in size from the published split, so it is unverified that the archive split keeps this. The value check also finds 208 test cases identical to training cases (Table~\ref{tab:duplicates}). Also, 54 of its repeated-value groups carry conflicting labels. Most matches are two-valued on/off days, so they need not mean shared households. \\
\hypertarget{split-audit:EthanolLevel}{}\SplitAuditName{EthanolLevel} & New bottle types & Separated & The archive paper says that a bottle type is never used in both training and testing \citep{dau2019ucr}. The test therefore checks measurements involving unfamiliar bottle types. This is a documented source separation. It does not, by itself, establish that every other recording condition is different between the sets. \\
\hypertarget{split-audit:FaceAll}{}\SplitAuditName{FaceAll} & New photography occasions & Concern & The task identifies people, so those identities are expected in both sets. The concern is whether test photographs were taken on a different occasion. The audit does not establish this. Similar photos taken under the same conditions may make identification easier than recognising the person in newly collected photographs. The value check also reports one test case identical to training data (Table~\ref{tab:duplicates}). \\
\hypertarget{split-audit:FaceFour}{}\SplitAuditName{FaceFour} & New source images or sessions & Unclear & The available record does not adequately describe the people, source photographs or capture sessions. We cannot check whether related images or acquisitions occur in both sets. No specific overlap is established by this absence of documentation. \\
\hypertarget{split-audit:FacesUCR}{}\SplitAuditName{FacesUCR} & New photography occasions & Unclear & The task uses reprocessed photographs from the same collection as FaceAll. Identifying familiar people is expected, but the audit does not establish that test photos were taken on a separate occasion. Similar lighting, pose or camera conditions could help recognition without showing that it works in newly collected photographs. \\
\hypertarget{split-audit:FiftyWords}{}\SplitAuditName{FiftyWords} & New manuscript pages & Unclear & The word images come from one writer's manuscripts. Words on the same page can share handwriting and scanning conditions. We do not know whether whole pages were kept in one set. The test may therefore contain more words from familiar pages, rather than words from pages the model has never seen. \\
\hypertarget{split-audit:Fish}{}\SplitAuditName{Fish} & New physical specimens & Unclear & The cases are outlines derived from fish images. The description does not establish how physical specimens and any repeated views were assigned to the two sets. New-specimen evaluation therefore cannot be verified; repeated views are a possibility to check, not an observed overlap. \\
\hypertarget{split-audit:FordA}{}\SplitAuditName{FordA} & New vehicles or recording occasions & Unclear & The files do not identify the vehicle, test rig or recording occasion behind each case. We therefore cannot check whether these sources were kept apart. The split may be appropriate, but the available information does not show that testing uses vehicles or recording occasions absent from training. \\
\hypertarget{split-audit:FordB}{}\SplitAuditName{FordB} & New conditions and recordings & Unclear & The audit describes test data collected under different, noisier operating conditions. This tests a change in conditions; it does not demonstrate shared recordings. We still do not know which vehicles or recording occasions produced the cases, so we cannot check whether those sources are different between the sets. \\
\hypertarget{split-audit:FreezerRegularTrain}{}\SplitAuditName{FreezerRegularTrain} & New physical freezers & Concern & The source contains ongoing measurements of two freezers in one house. Pieces of the same recording can share appliance settings and operating conditions. The audit does not show that testing was kept for a later period or a new recording occasion. The value check separately reports 1 of 2,850 test cases matching training data. \\
\hypertarget{split-audit:FreezerSmallTrain}{}\SplitAuditName{FreezerSmallTrain} & New physical freezers & Concern & Same as \hyperlink{split-audit:FreezerRegularTrain}{\SplitAuditName{FreezerRegularTrain}}. This smaller-training-set variant has 27 of 2,850 test cases matching training data in the value check. \\
\hypertarget{split-audit:Fungi}{}\SplitAuditName{Fungi} & New biological strains & Concern & The source measured fungal strains repeatedly. The audit reports that testing includes another measurement of a strain used in training. A model could recognise that familiar strain. This is a narrower test than identifying the species of a newly collected strain, for which all measurements of a strain should be kept together. \\
\hypertarget{split-audit:GestureMidAirD1}{}\SplitAuditName{GestureMidAirD1} & New people & Separated & The archive description says that no person occurs in both sets \citep{tscarchive_descriptions}. The test therefore uses people absent from training. There is disagreement in the records about which participant group is training, but both describe keeping the groups separate. No shared-person concern is identified by that description. \\
\hypertarget{split-audit:GestureMidAirD2}{}\SplitAuditName{GestureMidAirD2} & New people & Separated & Same as \hyperlink{split-audit:GestureMidAirD1}{\SplitAuditName{GestureMidAirD1}}. \\
\hypertarget{split-audit:GestureMidAirD3}{}\SplitAuditName{GestureMidAirD3} & New people & Separated & Same as \hyperlink{split-audit:GestureMidAirD1}{\SplitAuditName{GestureMidAirD1}}. \\
\hypertarget{split-audit:GesturePebbleZ1}{}\SplitAuditName{GesturePebbleZ1} & New people & Concern & The same four users performed gestures on two occasions. Training uses the first occasion and testing the second \citep{dau2019ucr}. Every test user is also in training, so the score can rest on per-user motion signatures and does not measure recognition for new wearers of the watch; \hyperlink{split-audit:GesturePebbleZ2}{\SplitAuditName{GesturePebbleZ2}} is the subject-disjoint version. \\
\hypertarget{split-audit:GesturePebbleZ2}{}\SplitAuditName{GesturePebbleZ2} & New people & Separated & Training uses two people and testing the other two \citep{dau2019ucr}. The model therefore sees unfamiliar users in the test. This differs from measuring familiar users again. The description establishes separation of people, although it does not establish separation of every other possible recording condition. \\
\hypertarget{split-audit:GunPoint}{}\SplitAuditName{GunPoint} & New people or recording occasions & Concern & All 200 series come from one recording session of two actors, so training and test necessarily share that session and the split cannot test newly collected recordings. Which actor is on which side is not recorded. Having only two actors also limits how broadly results can be interpreted. \\
\hypertarget{split-audit:GunPointAgeSpan}{}\SplitAuditName{GunPointAgeSpan} & New people or recording occasions & Concern & The archive pools actions from two actors recorded in two years \citep{dau2019ucr}. A check of the shipped files finds every actor and year combination in both sets, so the same two people and recording sessions appear on both sides. The score therefore concerns more actions from that collection. It does not establish performance on unfamiliar people or measurements collected on a new occasion. \\
\hypertarget{split-audit:GunPointMaleVersusFemale}{}\SplitAuditName{GunPointMaleVersusFemale} & New people or recording occasions & Concern & The two classes distinguish the two actors, and recordings from both years are pooled \citep{dau2019ucr}. A model can recognise these particular people to predict the label. The audit reports no new recording occasion in testing. The result also does not establish the same distinction for other people. \\
\hypertarget{split-audit:GunPointOldVersusYoung}{}\SplitAuditName{GunPointOldVersusYoung} & New recordings for each age condition & Concern & The label distinguishes the original recording years \citep{dau2019ucr}. Both sets use actions recorded on those occasions. Differences in equipment or recording conditions could help the model along with age-related changes. Further recordings under each age condition would be needed to check whether the result reflects age rather than those occasions. \\
\hypertarget{split-audit:Ham}{}\SplitAuditName{Ham} & New physical samples & Unclear & Each physical ham supplies several image measurements. We do not know whether all measurements of a ham were kept in one set. The test could therefore contain another measurement of a familiar ham. Not finding very similar pairs in the original check does not prove that different hams were used. \\
\hypertarget{split-audit:HandOutlines}{}\SplitAuditName{HandOutlines} & New patients & Unclear & The outlines come from hand X-rays. We do not know whether one patient contributed several images or whether those images were kept together. The test could contain a new image of a familiar patient. Patient identifiers or an explicit patient grouping rule are needed to check this. \\
\hypertarget{split-audit:Haptics}{}\SplitAuditName{Haptics} & New recording occasions & Unclear & The labels identify users from repeated measurements. We do not know whether the test measurements were collected on occasions separate from training. Familiar identities are expected, but the model might also recognise conditions specific to one recording. Identification when a user is measured again remains unverified. \\
\hypertarget{split-audit:Herring}{}\SplitAuditName{Herring} & New physical specimens & Unclear & The cases are outlines derived from fish ear-bone images. The description does not establish how physical specimens and any repeated views were assigned to the two sets. New-specimen evaluation therefore cannot be verified; repeated views are a possibility to check, not an observed overlap. \\
\hypertarget{split-audit:HouseTwenty}{}\SplitAuditName{HouseTwenty} & New households & Concern & The cases are pieces of one household's ongoing recording, so training and testing share the same home, its appliances and its routine. The test measures recognition within that home; performance in another home, which load monitoring needs, is not tested. \\
\hypertarget{split-audit:InlineSkate}{}\SplitAuditName{InlineSkate} & New recording occasions & Concern & Six of the 550 test series are value-identical to training series (Table~\ref{tab:duplicates}). If the classes are individuals, training and test cycles may also come from the same treadmill run and share its person, equipment and conditions, but no source confirms this. The test does not show whether recognition works after collecting another recording. This interpretation also needs confirmation of what the original class labels identify. \\
\hypertarget{split-audit:InsectEPGRegularTrain}{}\SplitAuditName{InsectEPGRegularTrain} & New insects or recordings & Unclear & The cases are pieces of long insect recordings. No insect or recording identifier ships with the data, so we cannot check whether pieces of one insect were kept in one set. Pieces in both sets could share the plant, sensor setup and nearby feeding behaviour, but this is not established. Unlike \hyperlink{split-audit:InsectEPGSmallTrain}{\SplitAuditName{InsectEPGSmallTrain}}, which shares its test set, this training set has no exact matches with the test data. \\
\hypertarget{split-audit:InsectEPGSmallTrain}{}\SplitAuditName{InsectEPGSmallTrain} & New insects or recordings & Concern & The cases are pieces of long insect recordings. No insect or recording identifier ships with the data, so whether pieces of one insect occur in both sets is not established. The value check finds 15 of the 249 test series identical to series in this small training set (Table~\ref{tab:duplicates}), so a model can score them by recall. \\
\hypertarget{split-audit:InsectWingbeatSound}{}\SplitAuditName{InsectWingbeatSound} & New insect groups or recordings & Unclear & The source recorded insects in groups inside cages. Cases from one group can share its insects and recording conditions. The files do not identify those groups, so we cannot check whether a whole group was kept in one set. Shared groups are possible, but have not been verified. \\
\hypertarget{split-audit:ItalyPowerDemand}{}\SplitAuditName{ItalyPowerDemand} & New days or later periods & Related periods & Training and test days can share the same season and similar weather. Different days are tested, so this is not automatically leakage. However, the split does not specifically test a later period or another grid. \\
\hypertarget{split-audit:LargeKitchenAppliances}{}\SplitAuditName{LargeKitchenAppliances} & New appliances or households & Unclear & Same as \hyperlink{split-audit:Computers}{\SplitAuditName{Computers}}. \\
\hypertarget{split-audit:Lightning2}{}\SplitAuditName{Lightning2} & New recorded events & Unclear & Each series is one recorded lightning event, but no source documents how the 121 records were divided or whether events of one flash or storm sit on both sides. \\
\hypertarget{split-audit:Lightning7}{}\SplitAuditName{Lightning7} & New recorded events & Unclear & Each series is one recorded lightning event. The source evaluated by random cross-validation over all 143 records \citep{eads2002zeus}; who built the fixed 70/73 split, and whether strokes of one flash sit on both sides, is not documented. \\
\hypertarget{split-audit:Mallat}{}\SplitAuditName{Mallat} & New generated examples & Generated & This is a generated task rather than measurements from people, devices or specimens. The corresponding source-separation question is therefore not assessed here. Its split would instead need to be examined through the generation procedure; synthetic origin alone is not evidence of leakage. \\
\hypertarget{split-audit:Meat}{}\SplitAuditName{Meat} & New physical samples & Unclear & The source measured each of 60 physical samples twice. We do not know whether the two measurements of a sample were kept together. If one goes into training and the other into testing, the model sees another measurement of a familiar sample. The sample assignment is needed to check whether this happened. \\
\hypertarget{split-audit:MedicalImages}{}\SplitAuditName{MedicalImages} & New patients or images & Concern & The files do not identify the patient or original image behind each series. Different series might therefore come from one patient or image. We cannot check whether those sources were kept apart. Missing identifiers leave this question unanswered; they do not prove that sources were shared. Separately, the value check reports one test case identical to training data (Table~\ref{tab:duplicates}). \\
\hypertarget{split-audit:MelbournePedestrian}{}\SplitAuditName{MelbournePedestrian} & New days or later periods & Related periods & Training and test days come from the same ten city sensors and the same year 2017; the assignment rule is undocumented. Training and test days can share local events, seasons and familiar sensors. The test checks more days from that collection. It does not specifically check a later period or a sensor that was absent from training. \\
\hypertarget{split-audit:MiddlePhalanxOutlineAgeGroup}{}\SplitAuditName{MiddlePhalanxOutlineAgeGroup} & New patients & Unclear & Same as \hyperlink{split-audit:DistalPhalanxOutlineAgeGroup}{\SplitAuditName{DistalPhalanxOutlineAgeGroup}}. \\
\hypertarget{split-audit:MiddlePhalanxOutlineCorrect}{}\SplitAuditName{MiddlePhalanxOutlineCorrect} & New patients & Unclear & Same as \hyperlink{split-audit:DistalPhalanxOutlineAgeGroup}{\SplitAuditName{DistalPhalanxOutlineAgeGroup}}. \\
\hypertarget{split-audit:MiddlePhalanxTW}{}\SplitAuditName{MiddlePhalanxTW} & New patients & Unclear & Same as \hyperlink{split-audit:DistalPhalanxOutlineAgeGroup}{\SplitAuditName{DistalPhalanxOutlineAgeGroup}}. \\
\hypertarget{split-audit:MixedShapesRegularTrain}{}\SplitAuditName{MixedShapesRegularTrain} & New images or physical objects & Concern & The cases are outlines from shape images. We do not know whether several views or processed versions of one original object can occur in both sets. If they do, the model may recognise a familiar shape. The source-image information needed to check this is not described in the audit. Separately, the value check reports one test case identical to training data (Table~\ref{tab:duplicates}). \\
\hypertarget{split-audit:MixedShapesSmallTrain}{}\SplitAuditName{MixedShapesSmallTrain} & New source images or objects & Unclear & Same uncertainty about the original images as \hyperlink{split-audit:MixedShapesRegularTrain}{\SplitAuditName{MixedShapesRegularTrain}}. The one value match mentioned there concerns MixedShapesRegularTrain, not this variant. \\
\hypertarget{split-audit:MoteStrain}{}\SplitAuditName{MoteStrain} & Independent recordings & Concern & The audit traces cases to nearby pieces of a small number of sensor recordings, with the recording defining the class. Training and testing therefore share those original recordings. The model may recognise their particular conditions. The test does not show that the distinction works across recordings collected independently. \\
\hypertarget{split-audit:NonInvasiveFetalECGThorax1}{}\SplitAuditName{NonInvasiveFetalECGThorax1} & Independent recordings & Unclear & The original audit thinks the labels identify recordings and that training and test heartbeats come from those same recordings. However, the source mapping is not confirmed. If correct, testing checks more beats from familiar recordings. This remains a possible concern until the original sources and labels are confirmed. \\
\hypertarget{split-audit:NonInvasiveFetalECGThorax2}{}\SplitAuditName{NonInvasiveFetalECGThorax2} & Independent source recordings & Unclear & Same as \hyperlink{split-audit:NonInvasiveFetalECGThorax1}{\SplitAuditName{NonInvasiveFetalECGThorax1}}. \\
\hypertarget{split-audit:OSULeaf}{}\SplitAuditName{OSULeaf} & New physical specimens & Unclear & The cases are outlines derived from leaf images. The description does not establish how physical specimens and any repeated views were assigned to the two sets. New-specimen evaluation therefore cannot be verified; repeated views are a possibility to check, not an observed overlap. \\
\hypertarget{split-audit:OliveOil}{}\SplitAuditName{OliveOil} & New physical samples & Unclear & The cases are spectral measurements of oils. The source measured 60 oils in duplicate; the archive ships 60 spectra and does not say whether they are one per oil or both replicates of 30 oils. Looking for similar sequences cannot tell us which sample produced them. The test's use of genuinely new samples therefore remains unverified. \\
\hypertarget{split-audit:PLAID}{}\SplitAuditName{PLAID} & New physical appliances & Concern & The source recorded individual appliances several times. We do not know whether all recordings of an appliance were kept in one set. Familiar appliances might occur in testing, but this is not established by the missing information. Separately, the value check finds five test cases identical to training cases. \\
\hypertarget{split-audit:PhalangesOutlinesCorrect}{}\SplitAuditName{PhalangesOutlinesCorrect} & New X-rays or patients & Unclear & Several bone outlines can come from one hand X-ray. Unless kept together, one image can supply outlines to both sets, letting the model benefit from familiar image conditions. The audit does not describe such grouping. Image or patient identifiers are needed to establish whether this possible overlap actually occurs. \\
\hypertarget{split-audit:Phoneme}{}\SplitAuditName{Phoneme} & New speakers or spoken recordings & Unclear & One spoken word can supply several phoneme segments. We do not know whether complete word recordings or speakers were kept in one set. Segments in both sets could share a voice and recording conditions. The test's use of new speakers or separately recorded words remains unverified. \\
\hypertarget{split-audit:PickupGestureWiimoteZ}{}\SplitAuditName{PickupGestureWiimoteZ} & New recording occasions & Unclear & The labels identify people, with several gesture repetitions per person. We do not know whether those repetitions were recorded on separate occasions or how the occasions were assigned to the sets. The test may check more repetitions from a familiar recording. Identification after recording the person again is not established. \\
\hypertarget{split-audit:PigAirwayPressure}{}\SplitAuditName{PigAirwayPressure} & New recording occasions & Concern & The task is to identify pigs. Training uses the first pieces of each pig's before/after pressure recordings; testing uses later pieces of the same recordings \citep{dau2019ucr}. The same pigs are expected in an identification task. The concern is that the model is tested on familiar recordings, rather than after measuring those pigs again. \\
\hypertarget{split-audit:PigArtPressure}{}\SplitAuditName{PigArtPressure} & New sessions of known animals & Concern & Same as \hyperlink{split-audit:PigAirwayPressure}{\SplitAuditName{PigAirwayPressure}}. \\
\hypertarget{split-audit:PigCVP}{}\SplitAuditName{PigCVP} & New sessions of known animals & Concern & Same as \hyperlink{split-audit:PigAirwayPressure}{\SplitAuditName{PigAirwayPressure}}. \\
\hypertarget{split-audit:Plane}{}\SplitAuditName{Plane} & New physical models or capture sessions & Unclear & The photographs show die-cast plane models, 30 shape samples per class; the source does not say how many physical models were photographed. The description does not establish how physical models or photography sessions were assigned to the two sets. The test could therefore contain another view of a familiar model; that possibility is not verified from the metadata. \\
\hypertarget{split-audit:PowerCons}{}\SplitAuditName{PowerCons} & New days or later periods & Concern & Nearby days can share weather and routines, and no source says whether a later period was kept for testing. Separately, 67 of 180 test cases exactly match training cases (Table~\ref{tab:duplicates}). These matches are a concrete issue beyond the possible relationship between neighbouring days. \\
\hypertarget{split-audit:ProximalPhalanxOutlineAgeGroup}{}\SplitAuditName{ProximalPhalanxOutlineAgeGroup} & New patients & Unclear & Same as \hyperlink{split-audit:DistalPhalanxOutlineAgeGroup}{\SplitAuditName{DistalPhalanxOutlineAgeGroup}}. \\
\hypertarget{split-audit:ProximalPhalanxOutlineCorrect}{}\SplitAuditName{ProximalPhalanxOutlineCorrect} & New patients & Unclear & Same as \hyperlink{split-audit:DistalPhalanxOutlineAgeGroup}{\SplitAuditName{DistalPhalanxOutlineAgeGroup}}. \\
\hypertarget{split-audit:ProximalPhalanxTW}{}\SplitAuditName{ProximalPhalanxTW} & New patients & Unclear & Same as \hyperlink{split-audit:DistalPhalanxOutlineAgeGroup}{\SplitAuditName{DistalPhalanxOutlineAgeGroup}}. \\
\hypertarget{split-audit:RefrigerationDevices}{}\SplitAuditName{RefrigerationDevices} & New appliances or households & Unclear & Same as \hyperlink{split-audit:Computers}{\SplitAuditName{Computers}}. \\
\hypertarget{split-audit:Rock}{}\SplitAuditName{Rock} & New physical samples & Unclear & The original check found strongly similar pairs of sequences. That alone does not show that the same physical rock was measured twice. We need sample identifiers to know whether the two sets share rocks. Until then, repeated measurements are a possibility, rather than confirmed overlap. \\
\hypertarget{split-audit:ScreenType}{}\SplitAuditName{ScreenType} & New appliances or households & Unclear & Same as \hyperlink{split-audit:Computers}{\SplitAuditName{Computers}}. \\
\hypertarget{split-audit:SemgHandGenderCh2}{}\SplitAuditName{SemgHandGenderCh2} & New people & Concern & The original audit reports the same five people in both sets, checked using labels from the related person-identification dataset. Recognising a familiar person could help predict their gender label. The test checks further measurements of these people, rather than establishing performance on people absent from training. \\
\hypertarget{split-audit:SemgHandMovementCh2}{}\SplitAuditName{SemgHandMovementCh2} & New people or recording occasions & Concern & The audit reports repeated grasps by the same people in the same recording sessions across both sets. These repetitions share the person and electrode placement. The test does not show whether recognition works for a new person or after removing and attaching the electrodes again. Those are separate evaluation settings. \\
\hypertarget{split-audit:SemgHandSubjectCh2}{}\SplitAuditName{SemgHandSubjectCh2} & New recording occasions & Concern & The task identifies people, but each person was recorded on one occasion according to the audit. Both sets contain repetitions from that occasion. The model might recognise that recording setup along with the person. The test does not establish identification after measuring the person again on another occasion. \\
\hypertarget{split-audit:ShakeGestureWiimoteZ}{}\SplitAuditName{ShakeGestureWiimoteZ} & New sessions of known users & Unclear & Same as \hyperlink{split-audit:PickupGestureWiimoteZ}{\SplitAuditName{PickupGestureWiimoteZ}}. \\
\hypertarget{split-audit:ShapeletSim}{}\SplitAuditName{ShapeletSim} & New generated examples & Generated & This is a generated task rather than measurements from people, devices or specimens. The corresponding source-separation question is therefore not assessed here. Its split would instead need to be examined through the generation procedure; synthetic origin alone is not evidence of leakage. \\
\hypertarget{split-audit:ShapesAll}{}\SplitAuditName{ShapesAll} & New source images or objects & Concern & The cases are outlines derived from shape images. The documentation does not establish whether repeated views or versions of one source object were kept in one set. Independence at the original image or object level remains unknown; no particular shared object is verified here. Separately, the value check reports two test cases identical to training data (Table~\ref{tab:duplicates}). \\
\hypertarget{split-audit:SmallKitchenAppliances}{}\SplitAuditName{SmallKitchenAppliances} & New appliances or households & Unclear & Same as \hyperlink{split-audit:Computers}{\SplitAuditName{Computers}}. \\
\hypertarget{split-audit:SmoothSubspace}{}\SplitAuditName{SmoothSubspace} & New generated examples & Generated & This is a generated task rather than measurements from people, devices or specimens. The corresponding source-separation question is therefore not assessed here. Its split would instead need to be examined through the generation procedure; synthetic origin alone is not evidence of leakage. \\
\hypertarget{split-audit:SonyAIBORobotSurface1}{}\SplitAuditName{SonyAIBORobotSurface1} & New recording occasions & Concern & The audit describes short pieces from one continuous robot recording per surface. Training and testing therefore share the robot run and its recording conditions. The test checks more pieces from those runs. It does not show whether surface recognition works after collecting another recording on a separate occasion. \\
\hypertarget{split-audit:SonyAIBORobotSurface2}{}\SplitAuditName{SonyAIBORobotSurface2} & New recording occasions & Concern & Same as \hyperlink{split-audit:SonyAIBORobotSurface1}{\SplitAuditName{SonyAIBORobotSurface1}}. \\
\hypertarget{split-audit:StarLightCurves}{}\SplitAuditName{StarLightCurves} & New stars & Separated & The source-level record describes one separate star per series. On that description, the test contains stars absent from training, rather than another piece of a familiar star's recording. No shared-star concern is identified. Possible shared observing conditions or processing procedures were not separately checked. \\
\hypertarget{split-audit:Strawberry}{}\SplitAuditName{Strawberry} & New physical samples & Concern & The test set is described as the source study's independently collected test data. However, the value check finds one sequence present in both sets with different labels (Table~\ref{tab:duplicates}). Thus, the original split description does not resolve this inconsistency in the released data. \\
\hypertarget{split-audit:SwedishLeaf}{}\SplitAuditName{SwedishLeaf} & New physical specimens & Concern & The cases are outlines derived from leaf images. The description does not establish how physical specimens and any repeated views were assigned to the two sets. New-specimen evaluation therefore cannot be verified; repeated views are a possibility to check, not an observed overlap. Separately, the value check reports two test cases identical to training data (Table~\ref{tab:duplicates}). \\
\hypertarget{split-audit:Symbols}{}\SplitAuditName{Symbols} & New people & Unclear & The source contains several symbol-drawing attempts per person. We do not know whether all attempts from a person were kept in one set. A test attempt could therefore come from a person already represented in training. This possible sharing needs participant identifiers to check; performance on new people is not established. \\
\hypertarget{split-audit:SyntheticControl}{}\SplitAuditName{SyntheticControl} & New generated examples & Generated & This is a generated task rather than measurements from people, devices or specimens. The corresponding source-separation question is therefore not assessed here. Its split would instead need to be examined through the generation procedure; synthetic origin alone is not evidence of leakage. \\
\hypertarget{split-audit:ToeSegmentation1}{}\SplitAuditName{ToeSegmentation1} & New people or recording occasions & Concern & The source contains repeated motion recordings from a small group of actors. The original audit suspects that actors are shared between the sets, but the files lack actor identifiers for checking this. We therefore cannot establish performance on new people or new recording occasions, and should not present actor overlap as verified. Separately, the value check finds four test cases identical to training data (Table~\ref{tab:duplicates}). \\
\hypertarget{split-audit:ToeSegmentation2}{}\SplitAuditName{ToeSegmentation2} & New people or recording sessions & Unclear & Same uncertainty about the actors and recording occasions as \hyperlink{split-audit:ToeSegmentation1}{\SplitAuditName{ToeSegmentation1}}. The four value matches mentioned there concern ToeSegmentation1, not this dataset. \\
\hypertarget{split-audit:Trace}{}\SplitAuditName{Trace} & New simulated events & Separated & Each series is described as a separate simulated event, so testing uses other generated events rather than another piece of one recording. No shared-event issue is identified by that description. Whether the generation process creates further dependencies would need its settings or code to check. \\
\hypertarget{split-audit:TwoLeadECG}{}\SplitAuditName{TwoLeadECG} & Independent recordings & Concern & The label identifies which ECG lead produced the signal. Training and test beats come from one patient's recording. The model may recognise differences specific to those leads during that measurement. The test does not show that the same distinction works when another patient or recording is used. \\
\hypertarget{split-audit:TwoPatterns}{}\SplitAuditName{TwoPatterns} & New generated examples & Generated & This is a generated task rather than measurements from people, devices or specimens. The corresponding source-separation question is therefore not assessed here. Its split would instead need to be examined through the generation procedure; synthetic origin alone is not evidence of leakage. \\
\hypertarget{split-audit:UMD}{}\SplitAuditName{UMD} & New generated examples & Generated & This is a generated task rather than measurements from people, devices or specimens. The corresponding source-separation question is therefore not assessed here. Its split would instead need to be examined through the generation procedure; synthetic origin alone is not evidence of leakage. \\
\hypertarget{split-audit:UWaveGestureLibraryAll}{}\SplitAuditName{UWaveGestureLibraryAll} & New days or recording occasions & Unclear & The original system recognises gestures for familiar users, so sharing users can be appropriate. The missing information is whether all recordings from a day were kept together. Without this, the test could contain gestures recorded on a day already represented in training, rather than gestures collected when the user returns. \\
\hypertarget{split-audit:UWaveGestureLibraryX}{}\SplitAuditName{UWaveGestureLibraryX} & New sessions or days & Unclear & Same as \hyperlink{split-audit:UWaveGestureLibraryAll}{\SplitAuditName{UWaveGestureLibraryAll}}. \\
\hypertarget{split-audit:UWaveGestureLibraryY}{}\SplitAuditName{UWaveGestureLibraryY} & New sessions or days & Unclear & Same as \hyperlink{split-audit:UWaveGestureLibraryAll}{\SplitAuditName{UWaveGestureLibraryAll}}. \\
\hypertarget{split-audit:UWaveGestureLibraryZ}{}\SplitAuditName{UWaveGestureLibraryZ} & New sessions or days & Unclear & Same as \hyperlink{split-audit:UWaveGestureLibraryAll}{\SplitAuditName{UWaveGestureLibraryAll}}. \\
\hypertarget{split-audit:Wafer}{}\SplitAuditName{Wafer} & New wafers or production runs & Concern & The archive ships the six sensor traces of each wafer as separate cases and splits them without regard to wafer or run. Matching every case to the raw files shows 690 wafers and 99 of 104 production runs on both sides; every training case has another sensor of its wafer in test. The test therefore does not measure performance on new wafers or runs. \\
\hypertarget{split-audit:Wine}{}\SplitAuditName{Wine} & New physical samples & Unclear & We do not know how many physical wines produced the spectral measurements or whether one wine was measured repeatedly. The two sets might contain measurements of the same wine. Not finding obvious similar pairs does not settle this question. Physical-sample identifiers are needed to check the split. \\
\hypertarget{split-audit:WordSynonyms}{}\SplitAuditName{WordSynonyms} & New manuscript pages & Unclear & Same as \hyperlink{split-audit:FiftyWords}{\SplitAuditName{FiftyWords}}. \\
\hypertarget{split-audit:Worms}{}\SplitAuditName{Worms} & New worms or recordings & Concern & The original record describes each case as a separate worm recording. However, the value check finds 10 of 77 test cases identical to training cases (Table~\ref{tab:duplicates}). We therefore need worm and recording identifiers to resolve the disagreement between that source description and the released values. \\
\hypertarget{split-audit:WormsTwoClass}{}\SplitAuditName{WormsTwoClass} & New worms or acquisitions & Concern & Same as \hyperlink{split-audit:Worms}{\SplitAuditName{Worms}}. \\
\hypertarget{split-audit:Yoga}{}\SplitAuditName{Yoga} & New filming occasions & Concern & The task identifies two actors from video frames. No test frame is identical to a training frame, but every training frame has a near-twin in test (median distance 2.09 versus 23.8 for random pairs), so both sets come from one pool of frames. The model may recognise familiar filming conditions along with the actor. The test does not establish identification in footage collected on another occasion. \\
\end{longtable}
\endgroup
